% v10 common-socle, full-record vs binary trace bridge; 9 October 2026.
% Candidate replacement of the prior Lemma 6.49 TODO in the cumulative
% author-reviewed v10 source. Do not insert duplicates or edit the
% author's four reconstructed mathematical proof replacements.
%
\todo[inline]{\textbf{v10 validation (first complete type difference).}
For two partial types extracted from the same ambient $C^+$,
assume the level-zero types agree and restrict both below their
free levels. Then the entire induced $L^+$-types coincide through
$d$ if and only if both directed binary relation vectors between
their distinguished type vertices and every common socle coordinate
below $d$ coincide. The ambient socle supplies the old
$L^+$-relations, equality of the roots supplies the singleton and
diagonal data (including loops), and the free-level cut determines
every auxiliary $E$-pair incident with the type vertex. The
implication in both directions, and the passage from the first
binary disagreement at $d$ to equality of the complete predecessors
through $d$ but not $d+1$, are proved in
\texttt{SuccessorTree/V10/CommonSocleType.lean}, conditional on the
two restrictions being defined, notably $d+1\leq\min
\{\mathrm{fl}_{C^+}(v),\mathrm{fl}_{C^+}(w)\}$.
For Lemma~6.49 still show that the proposed positive-generation
coordinate satisfies this cut restriction (using the sharp
$p+1<j$ condition where applicable) and identify these complete
prefix records with the actual $K^\mathcal F_{pt}$ tree
ancestors before invoking the checked generic
\texttt{RecordBridge.lean} meet result.}
%
% Endpoints, subject to the exact SHA's successful CI:
%   sameAmbient_partialType_eq_of_cross
%   sameAmbient_cross_of_partialType_eq
%   sameAmbient_partialType_eq_iff_cross
%   sameAmbient_first_fullType_difference
